\section{Strong universality: registers that build themselves}
\label{sec:strong}

This is the central section. The motion and switch rules of Section~\ref{sec:automaton} already give a universal
automaton on every tiling, but only weakly, because the unbounded registers of a register machine have to be
supplied as an infinite, if ultimately periodic, initial track. We now remove that infinity. We add one rule,
a track-building rule, under which a register starts as a single cell and constructs the digits it needs as the
computation runs. The program is finite and the registers self-extend, so the whole machine starts from a
finite configuration, which is strong universality.

\subsection{The track-building rule}

Empty space is not inert at the growing tip of a register. When the locomotive reaches the end of a register
and must extend it, the head selects a fresh empty neighbour, turns it into track, and advances into it. The
rule is Table~\ref{tab:build}. To choose the fresh cell with no global frame, the head takes the empty
neighbour of greatest tree depth, the outward direction laid down with the seed, breaking ties by the cell's
own index. On a tiling whose balls grow without bound a strictly deeper empty cell always exists, so the
builder never traps itself.

\begin{table}[ht]
\centering
\begin{tabular}{@{}lp{10.0cm}@{}}
\toprule
rule & action \\
\midrule
(B1) & a head on the end of a register, with no clear track ahead, turns its deepest empty neighbour into a
track cell and advances the head into it \\
(B2) & the new track cell links back to the head's cell and forward to its own next outward cell, so the
register has grown by one digit \\
\bottomrule
\end{tabular}
\caption{The track-building rule, the single addition that turns weak universality into strong universality.
Outward is read from tree depth, which increases without bound, so the rule never traps.}
\label{tab:build}
\end{table}

\begin{lemma}[The builder never traps]
\label{lem:build}
On an infinite regular tiling, a builder governed by (B1) and (B2) takes a strictly outward step at every step,
so it constructs an unbounded track from a one-cell seed.
\end{lemma}

Every cell of an infinite regular tiling has a child of strictly greater tree depth, which is empty until the
builder reaches it, so the deepest-empty-neighbour choice is always defined and strictly increases the depth.
The track therefore grows by one cell at every step and never halts. \qed

We verified this directly. From a one-cell seed the builder lays a strictly outward track on the octagrid, the
dodecagrid, and the $\{3,4,3,4\}$ honeycomb, halting in a finite test patch only at the patch boundary and so
never on the infinite tiling \cite{pollard2026vibe}.

\subsection{The self-extending register}

A register is a chain of flip-flop bits holding a number in binary, the lowest bit at the entry. The locomotive
increments it by rippling, it flips the lowest bit, and on a carry flips the next, and so on. Table~\ref{tab:reg}
gives the operation. The new ingredient is the last line, on overflow, when the carry runs past the top bit,
the builder rule (B1) lays a new bit and the carry sets it. So the register that began with one bit grows a new
bit exactly when it first needs one.

\begin{table}[ht]
\centering
\begin{tabular}{@{}llp{7.4cm}@{}}
\toprule
the locomotive meets a bit & whose value is & and \\
\midrule
flip-flop bit, active crossing & $0$ & sets it to $1$, no carry, returns (increment done) \\
flip-flop bit, active crossing & $1$ & sets it to $0$, carries to the next bit \\
end of register, carry out & --- & builds a new bit by (B1) and sets it to $1$ (the register grows) \\
\bottomrule
\end{tabular}
\caption{The self-extending register. Increment ripples carries through the flip-flop bits, and on overflow the
register builds itself a new bit. The bits' selection settings are the stored number.}
\label{tab:reg}
\end{table}

\begin{lemma}[Self-extension]
\label{lem:selfext}
A register started as a single bit, incremented $k$ times, holds the value $k$ in $\lceil \log_2(k+1) \rceil$
bits, all of which past the first it built itself.
\end{lemma}

Each increment adds one by the binary ripple of Table~\ref{tab:reg}, and a new bit is built precisely when the
count first reaches a new power of two, by (B1). So after $k$ increments the bits hold $k$ and number
$\lceil \log_2(k+1) \rceil$. \qed

We verified this directly. A counter started from a one-bit seed counts $1$ through $100$ correctly, growing
from one bit to seven by six self-builds \cite{pollard2026vibe}. The growth on demand is the same move a growing
mesh makes when it adds a new ring at its open edge each step, memory is not laid out in advance, it is built as
the process reaches for it.

\subsection{The strong universality theorem}

\begin{theorem}[Uniform strong universality]
\label{thm:strong}
There is one cellular automaton, defined by rules local in the cell-adjacency graph and independent of any
tiling (Tables~\ref{tab:states} to~\ref{tab:reg}), such that for every regular hyperbolic tiling and every
register machine there is a \textit{finite} initial configuration on that tiling whose evolution simulates the machine.
\end{theorem}

The initial configuration holds the finite track of the machine's fixed program, the switches that wire its
instructions, and each register as a single seed bit. The locomotive runs the program by the railway reduction
of Proposition~\ref{prop:railway}, using the motion and switch rules of Section~\ref{sec:automaton}. Whenever an
increment overflows a register, the register extends itself by the build rule (B1), by Lemmas~\ref{lem:build}
and~\ref{lem:selfext}, so the unbounded memory the machine needs is constructed from the finite seed rather than
given. A register machine with two counters is universal \cite{minsky1967}, so the configuration is finite and
the evolution is universal, which is strong universality. None of the rules referred to the tiling, so one
automaton serves every regular tiling. \qed

\subsection{Scope, stated honestly}

The pieces of Theorem~\ref{thm:strong} are each verified by direct execution, the locomotive and the three
switches on the actual tilings (Section~\ref{sec:verification}), a register computed end to end as a binary
counter, the self-extending counter from a finite seed, and the never-trapping builder. The theorem composes
them through the standard railway-to-register-machine reduction, the same way Margenstern's strongly universal
dodecagrid automaton is established, by a construction together with checks of its parts
\cite{margenstern2021dodeca5}. We do not exhibit a single monolithic run of a complete two-register program
laid on one tiling from a finite seed, which is an engineering assembly of the verified parts rather than a new
idea, and we mark it as such.
